% ==============================================================================
% Figura: Arquitectura Global de Conmutación e Interconexión de Bloques (1 kHz)
% Estrategia: Layout by Construction, paleta e idioma visual idénticos a Figs. 7-10.
% Regla estricta: SOLO nombres cortos de bloques (cero subtítulos ni texto extra),
% realimentación separada por sensor (brazo theta vs péndulo alpha) y separada
% entre posición y velocidad (6 pistas ordenadas con puentes limpios).
% Bibliotecas TikZ requeridas: arrows.meta, positioning, fit, calc,
%                              shapes.geometric, backgrounds, babel, bending.
% ==============================================================================
\begin{figure*}[!t]
\centering
% --- Paleta semántica armónica (idéntica a Figs. 7, 8, 9 y 10) ---
\definecolor{blkInk}{HTML}{1E293B}
\definecolor{blkPendEdge}{HTML}{1B4965}
\definecolor{blkPendFill}{HTML}{EAF4FB}
\definecolor{blkArmEdge}{HTML}{0F766E}
\definecolor{blkArmFill}{HTML}{E6F7F5}
\definecolor{blkCondEdge}{HTML}{B45309}
\definecolor{blkCondFill}{HTML}{FEF6EB}
\definecolor{blkDistEdge}{HTML}{7E2268}
\definecolor{blkDistFill}{HTML}{FBF0F8}
\definecolor{blkActEdge}{HTML}{A63429}
\definecolor{blkActFill}{HTML}{FDF0EE}
\definecolor{blkPlantEdge}{HTML}{243447}
\definecolor{blkPlantFill}{HTML}{ECEFF4}
\definecolor{blkZoneZ1}{HTML}{33658A}
\definecolor{blkZoneL2}{HTML}{A63429}
\definecolor{blkZoneL3}{HTML}{2A6B43}
\definecolor{blkZoneSup}{HTML}{0F766E}

% --- Parámetros centralizados de ritmo geométrico (Single Source of Truth) ---
\def\cbModW{3.45cm}
\def\cbRowGap{0.34cm}
\def\cbMuxGap{1.35cm}
\def\cbMuxW{1.55cm}
\def\cbActGap{0.44cm}
\def\cbActW{1.82cm}
\def\cbActH{0.86cm}
\def\cbSensGap{0.46cm}
\def\cbSensW{2.04cm}
\def\cbSensH{0.82cm}
\def\cbTrkStep{0.46cm}
\def\cbHopRad{0.085cm}
\def\cbPortGap{0.40cm}

\resizebox{\textwidth}{!}{%
\begin{tikzpicture}[
    >=Stealth,
    font=\sffamily,
    % --- Estilos de bloques atómicos (solo nombre corto, sin texto extra) ---
    baseblk/.style={
        rectangle,
        rounded corners=3.5pt,
        align=center,
        inner xsep=3.5pt,
        inner ysep=3pt,
        text=blkInk,
        font=\small\sffamily\bfseries
    },
    modIdle/.style={baseblk, draw=blkPlantEdge, fill=blkPlantFill, line width=1.05pt, text width=\cbModW, minimum height=0.72cm},
    modArm/.style={baseblk, draw=blkArmEdge, fill=blkArmFill, line width=1.1pt, text width=\cbModW},
    modCond/.style={baseblk, draw=blkCondEdge, fill=blkCondFill, line width=1.1pt, text width=\cbModW},
    modAct/.style={baseblk, draw=blkActEdge, fill=blkActFill, line width=1.1pt, text width=\cbModW},
    modPend/.style={baseblk, draw=blkPendEdge, fill=blkPendFill, line width=1.45pt, text width=\cbModW},
    actblk/.style={baseblk, draw=blkActEdge, fill=blkActFill, line width=1.05pt, text width=\cbActW, minimum height=\cbActH},
    plantblk/.style={baseblk, draw=blkPlantEdge, fill=blkPlantFill, line width=1.35pt, text width=1.74cm, minimum height=\cbActH},
    sensTblk/.style={baseblk, draw=blkArmEdge, fill=blkArmFill, line width=1.05pt, text width=\cbSensW, minimum height=\cbSensH},
    sensAblk/.style={baseblk, draw=blkPendEdge, fill=blkPendFill, line width=1.05pt, text width=\cbSensW, minimum height=\cbSensH},
    suppill/.style={
        rectangle,
        rounded corners=4pt,
        align=center,
        minimum height=0.60cm,
        inner xsep=6pt,
        inner ysep=2.5pt,
        text=blkInk,
        font=\footnotesize\sffamily\bfseries,
        line width=1.05pt
    },
    % --- Contenedores y etiquetas (idénticos a Figs. 7, 8 y 10) ---
    zonebox/.style={
        rectangle,
        rounded corners=5pt,
        draw=#1,
        densely dashed,
        line width=0.95pt,
        fill=#1,
        fill opacity=0.035,
        text opacity=1
    },
    zonetitle/.style={
        rectangle,
        rounded corners=3pt,
        draw=#1,
        line width=0.8pt,
        fill=white,
        inner xsep=5.5pt,
        inner ysep=2.2pt,
        font=\small\sffamily\bfseries,
        text=#1
    },
    dot/.style={circle, fill=#1, inner sep=1.45pt},
    % --- Flechas canónicas ---
    sig/.style={-{Stealth[length=4.6pt, width=3.4pt]}, line width=1.0pt, draw=blkInk},
    sigshort/.style={-{Stealth[length=3.9pt, width=2.9pt]}, line width=0.95pt, draw=blkInk},
    fbT/.style={-{Stealth[length=4.2pt, width=3.0pt]}, line width=0.95pt, draw=blkArmEdge},
    fbA/.style={-{Stealth[length=4.2pt, width=3.0pt]}, line width=0.95pt, draw=blkPendEdge},
    lineT/.style={line width=0.95pt, draw=blkArmEdge},
    lineA/.style={line width=0.95pt, draw=blkPendEdge},
    adapt/.style={-{Stealth[length=4.4pt, width=3.2pt]}, line width=1.05pt, densely dashed, draw=blkDistEdge},
    byp/.style={-{Stealth[length=4.4pt, width=3.2pt]}, line width=1.0pt, densely dashed, draw=blkActEdge},
    wirelbl/.style={font=\footnotesize\sffamily, text=blkInk, inner sep=1.2pt}
]

% ==============================================================================
% 1. BANCO DE DIAGRAMAS DE BLOQUES CONMUTABLES (Solo nombre corto por bloque)
% ==============================================================================
\node[modIdle] (modIdle) at (0,0) {IDLE / FAULT};
\node[modArm, minimum height=0.80cm, below=\cbRowGap of modIdle] (modVel) {MOTOR\_PID\_VEL};
\node[modArm, minimum height=0.98cm, below=\cbRowGap of modVel] (modPos) {Servo $0^\circ$};
\node[modAct, minimum height=1.32cm, below=\cbRowGap of modPos] (modRec) {RECOVER\_HOME};
\node[modCond, minimum height=1.68cm, below=\cbRowGap of modRec] (modSw) {SWING\_UP};
\node[modPend, minimum height=2.08cm, below=\cbRowGap of modSw] (modBal) {BALANCE\_PID};

% Reutilización directa F1: RECOVER_HOME invoca Servo 0°
\draw[adapt, draw=blkArmEdge] (modRec.north) -- node[wirelbl, right=2.5pt, text=blkArmEdge]{\textbf{F1}} (modPos.south);

% ==============================================================================
% 2. CONMUTADOR CENTRAL MUX FSM (Construido sobre el rango vertical exacto)
% ==============================================================================
\coordinate (muxNW) at ($(modIdle.north east)+(\cbMuxGap, 0.14)$);
\coordinate (muxSW) at ($(modBal.south east)+(\cbMuxGap, -0.14)$);
\coordinate (muxNE) at ($(muxNW)+(\cbMuxW, -0.72)$);
\coordinate (muxSE) at ($(muxSW)+(\cbMuxW, 0.72)$);
\coordinate (muxMidR) at ($(muxNE)!0.5!(muxSE)$);
\coordinate (muxTop)  at ($(muxNW)!0.48!(muxNE)$);
\coordinate (muxOutU) at ($(muxNE)!0.38!(muxSE)$);
\coordinate (muxOutB) at ($(muxNE)!0.66!(muxSE)$);

\filldraw[fill=blkDistFill, draw=blkDistEdge, line width=1.35pt, rounded corners=4pt]
    (muxNW) -- (muxNE) -- (muxSE) -- (muxSW) -- cycle;
\node[font=\small\sffamily\bfseries, text=blkDistEdge, align=center]
    at ($(muxNW)!0.52!(muxNE)+(0,-0.56)$) {MUX\\[-2pt]FSM};

% Contactos y brazo selector rotatorio alineados a cada bloque
\coordinate (muxPivot) at ($(muxMidR)+(-0.34, 0)$);
\foreach \m/\c in {modIdle/blkPlantEdge, modVel/blkArmEdge, modPos/blkArmEdge, modRec/blkActEdge, modSw/blkCondEdge, modBal/blkPendEdge} {
    \coordinate (pIn-\m) at (muxNW |- \m.east);
    \coordinate (pDot-\m) at ($(pIn-\m)+(0.25, 0)$);
    \draw[draw=\c, line width=0.8pt] (pIn-\m) -- (pDot-\m);
    \fill[fill=\c] (pDot-\m) circle (1.9pt);
}
\draw[draw=blkDistEdge, line width=1.5pt, line cap=round] (muxPivot) -- (pDot-modBal);
\fill[fill=blkDistEdge] (muxPivot) circle (2.3pt);
\draw[draw=blkDistEdge!55, line width=0.85pt, densely dotted, -{Stealth[length=3.6pt,width=2.6pt]}]
    ($(muxPivot)+(-0.56, 1.75)$) to[bend right=18] ($(muxPivot)+(-0.56, -1.75)$);

% Salidas de bloques -> MUX (solo símbolos puros)
\draw[sigshort, draw=blkPlantEdge] (modIdle.east) -- node[wirelbl, pos=0.60, above=1.5pt]{$0$} (pIn-modIdle);
\draw[sigshort, draw=blkArmEdge]   (modVel.east)  -- node[wirelbl, pos=0.60, above=1.5pt]{$u_v$} (pIn-modVel);
\draw[sigshort, draw=blkArmEdge]   (modPos.east)  -- node[wirelbl, pos=0.60, above=1.5pt]{$u_{\text{servo}}$} (pIn-modPos);
\draw[sigshort, draw=blkActEdge]   (modRec.east)  -- node[wirelbl, pos=0.60, above=1.5pt]{\texttt{F0/F2}} (pIn-modRec);
\draw[sigshort, draw=blkCondEdge]  (modSw.east)   -- node[wirelbl, pos=0.60, above=1.5pt]{$u_{\text{sw}}$} (pIn-modSw);
\draw[sigshort, draw=blkPendEdge]  (modBal.east)  -- node[wirelbl, pos=0.60, above=1.5pt]{$u_{\text{bal}}$} (pIn-modBal);

% ==============================================================================
% 3. ACTUACIÓN Y PLANTA (CASCADA IDÉNTICA A FIGS. 7, 8 Y 10)
% ==============================================================================
\node[actblk, right=1.18cm of muxOutU] (banda) {Banda\\Silenciosa};
\node[actblk, right=\cbActGap of banda] (rampa) {Rampa\\L298N};
\node[actblk, right=\cbActGap of rampa] (pwm) {Modulador\\PWM};
\node[actblk, right=\cbActGap of pwm] (puente) {Puente H\\L298N};
\node[actblk, right=0.58cm of puente] (motor) {Motor\\80:1};
\node[plantblk, right=0.58cm of motor] (planta) {P\'endulo\\Furuta};

\draw[sig] (muxOutU) -- node[wirelbl, pos=0.36, above=1.5pt]{$u[k]$} (banda.west);
\draw[sigshort] (banda.east) -- (rampa.west);
\draw[sigshort] (rampa.east) -- (pwm.west);
\draw[sigshort] (pwm.east) -- (puente.west);
\draw[sigshort] (puente.east) -- node[wirelbl, midway, above=1.5pt]{$V_m$} (motor.west);
\draw[sigshort] (motor.east) -- node[wirelbl, midway, above=1.5pt]{$\tau_0$} (planta.west);

% Desvío directo Freno / Libre hacia Puente H L298N
\coordinate (bypTurn) at ($(banda.south west)+(0.10, -0.78)$);
\draw[byp, rounded corners=3pt]
    (muxOutB) -- (muxOutB -| bypTurn) |- (bypTurn -| puente.south) -- (puente.south);
\node[wirelbl, above=1.5pt, text=blkActEdge] (lblByp)
    at ($(bypTurn)!0.50!(bypTurn -| puente.south)$) {Freno / Libre};

% ==============================================================================
% 4. REALIMENTACIÓN Y FILTRADO PARALELO (IDÉNTICO A FIGS. 7 Y 8)
% ==============================================================================
\coordinate (sensYTop) at ($(modBal.south)+(0, -2.42)$);
\node[sensTblk, anchor=center] (as5600_base) at (motor.center |- sensYTop) {AS5600\\Base};
\node[sensTblk, left=\cbSensGap of as5600_base] (desenv) {Desenvoltura};
\node[sensTblk, left=\cbSensGap of desenv] (deriv_t) {Derivador $\theta$};
\node[sensTblk, left=\cbSensGap of deriv_t] (lpf1_t) {LPF IIR 1\\($\theta$)};
\node[sensTblk, left=\cbSensGap of lpf1_t] (lpf2_t) {LPF IIR 2\\($\theta$)};

\node[sensAblk, below=0.52cm of as5600_base] (as5600_dist) {AS5600\\Distal};
\node[sensAblk, left=\cbSensGap of as5600_dist] (envolt) {Envoltura};
\node[sensAblk, left=\cbSensGap of envolt] (deriv_a) {Derivador $\alpha$};
\node[sensAblk, left=\cbSensGap of deriv_a] (lpf1_a) {LPF IIR 1\\($\alpha$)};
\node[sensAblk, left=\cbSensGap of lpf1_a] (lpf2_a) {LPF IIR 2\\($\alpha$)};

% Dos salidas físicas separadas de la planta: theta(t) y alpha(t)
\coordinate (pl_t_out) at ([xshift=-0.26cm]planta.south);
\coordinate (pl_a_out) at ([xshift= 0.26cm]planta.south);
\draw[fbT, rounded corners=4pt] (pl_t_out) |- (as5600_base.east);
\node[wirelbl, text=blkArmEdge, anchor=east]
    at ([xshift=-0.08cm]$(pl_t_out |- banda.south)!0.65!(pl_t_out |- as5600_base.east)$) {$\theta(t)$};
\draw[fbA, rounded corners=4pt] (pl_a_out) |- (as5600_dist.east);
\node[wirelbl, text=blkPendEdge, anchor=west]
    at ([xshift=0.08cm]$(pl_a_out |- as5600_base.east)!0.50!(pl_a_out |- as5600_dist.east)$) {$\alpha(t)$};

\draw[fbT] (as5600_base.west) -- (desenv.east);
\draw[fbT] (desenv.west) -- (deriv_t.east);
\draw[fbT] (deriv_t.west) -- (lpf1_t.east);
\draw[fbT] (lpf1_t.west) -- (lpf2_t.east);

\draw[fbA] (as5600_dist.west) -- (envolt.east);
\draw[fbA] (envolt.west) -- (deriv_a.east);
\draw[fbA] (deriv_a.west) -- (lpf1_a.east);
\draw[fbA] (lpf1_a.west) -- (lpf2_a.east);

% ==============================================================================
% 5. CAPA SUPERIOR: SUPERVISOR FSM (SOLO NOMBRES DE ESTADOS Y TRANSICIONES)
% ==============================================================================
\coordinate (supMidX) at ($(modIdle.west)!0.42!(planta.east)$);
\coordinate (supY)    at ($(modIdle.north)+(0, 3.35)$);

\node[suppill, draw=blkCondEdge, fill=blkCondFill] (supSw) at ($(supMidX |- supY)+(-2.25, 0)$) {SWING\_UP};
\node[suppill, draw=blkPlantEdge, fill=blkPlantFill, left=1.10cm of supSw] (supIdle) {IDLE};
\node[suppill, draw=blkPendEdge, fill=blkPendFill, line width=1.4pt, right=1.15cm of supSw] (supBal) {BALANCE\_PID};
\node[suppill, draw=blkActEdge, fill=blkActFill, right=1.15cm of supBal] (supRec) {RECOVER\_HOME};
\node[suppill, draw=blkArmEdge, fill=blkArmFill, right=1.00cm of supRec] (supAux) {MOTOR\_PID / STOP\_HOME};
\node[suppill, draw=blkPlantEdge, fill=white, left=1.45cm of supIdle] (supCli) {CLI / B1};

\draw[sigshort] (supCli.east) -- node[wirelbl, pos=0.32, above=1.5pt]{cmd} (supIdle.west);
\draw[sigshort] (supIdle.east) -- node[wirelbl, midway, above=1.5pt]{\texttt{START}} (supSw.west);
\draw[sigshort, draw=blkPendEdge] (supSw.east) -- node[wirelbl, midway, above=1.5pt]{Captura} (supBal.west);
\draw[sigshort, draw=blkActEdge] (supBal.east) -- node[wirelbl, midway, above=1.5pt]{Ca\'ida} (supRec.west);

\coordinate (calmTop) at ($(supRec.north)+(0, 0.38)$);
\coordinate (calmLblPad) at ($(calmTop)+(0, 0.28)$);
\draw[sigshort, draw=blkCondEdge, rounded corners=3pt]
    (supRec.north) -- (calmTop) -| node[wirelbl, pos=0.25, above=1.5pt, text=blkCondEdge]{Calma} (supSw.north);

\coordinate (auxBot) at ($(supIdle.south)+(0, -0.34)$);
\draw[sigshort, draw=blkArmEdge, densely dashed, rounded corners=3pt]
    (supIdle.south) -- (auxBot) -| node[wirelbl, pos=0.25, below=1.5pt, text=blkArmEdge]{\texttt{VEL / POS / STOPHOME}} (supAux.south);
\coordinate (auxLblPad) at ($(auxBot)+(0, -0.32)$);

% ==============================================================================
% 6. CONTENEDORES SEMÁNTICOS POR CONSTRUCCIÓN (fit en capa de fondo)
% ==============================================================================
\begin{scope}[on background layer]
    \node[zonebox=blkZoneSup,
          fit=(supIdle)(supSw)(supBal)(supRec)(supAux)(calmLblPad)(auxLblPad),
          inner ysep=9pt, inner xsep=10pt] (zSup) {};
    \node[zonetitle=blkZoneSup, anchor=west]
          at ([xshift=0.35cm]zSup.north west) {Supervisor FSM};

    \node[zonebox=blkZoneZ1,
          fit=(modIdle)(modVel)(modPos)(modRec)(modSw)(modBal),
          inner ysep=10pt, inner xsep=9pt] (zBank) {};
    \node[zonetitle=blkZoneZ1, anchor=west]
          at ([xshift=0.30cm]zBank.north west) {Controladores por Estado};

    \node[zonebox=blkZoneL2,
          fit=(banda)(rampa)(pwm)(puente)(motor)(planta)(bypTurn)(lblByp),
          inner ysep=10pt, inner xsep=9pt] (zAct) {};
    \node[zonetitle=blkZoneL2, anchor=west]
          at ([xshift=0.32cm]zAct.north west) {Actuaci\'on y Planta (Cascada)};

    \node[zonebox=blkZoneL3,
          fit=(as5600_base)(desenv)(deriv_t)(lpf1_t)(lpf2_t)(as5600_dist)(envolt)(deriv_a)(lpf1_a)(lpf2_a),
          inner ysep=11pt, inner xsep=9pt] (zSens) {};
    \node[zonetitle=blkZoneL3, anchor=east]
          at ([xshift=-0.16cm]zSens.south east) {Realimentaci\'on y Filtrado};
\end{scope}

% ==============================================================================
% 7. 6 PISTAS DE REALIMENTACIÓN SEPARADAS (POSICIÓN != VELOCIDAD, THETA != ALPHA)
% ==============================================================================
% Coordenadas X de las 6 pistas verticales en el margen izquierdo (interior -> exterior)
\coordinate (trk_dt2) at ([xshift=-0.78cm]zBank.west);
\coordinate (trk_da2) at ([xshift=-\cbTrkStep]trk_dt2);
\coordinate (trk_dt1) at ([xshift=-\cbTrkStep]trk_da2);
\coordinate (trk_t)   at ([xshift=-\cbTrkStep]trk_dt1);
\coordinate (trk_da1) at ([xshift=-\cbTrkStep]trk_t);
\coordinate (trk_a)   at ([xshift=-\cbTrkStep]trk_da1);

% Puntos de derivación separados en zSens (posición tras Desenvoltura/Envoltura; velocidad tras LPF 1/2)
\coordinate (tap_t)   at ($(desenv.west)!0.52!(deriv_t.east)$);
\coordinate (tap_dt1) at ($(lpf1_t.west)!0.52!(lpf2_t.east)$);
\coordinate (tap_a)   at ($(envolt.west)!0.52!(deriv_a.east)$);
\coordinate (tap_da1) at ($(lpf1_a.west)!0.52!(lpf2_a.east)$);

\node[dot=blkArmEdge]  at (tap_t) {};
\node[dot=blkArmEdge]  at (tap_dt1) {};
\node[dot=blkPendEdge] at (tap_a) {};
\node[dot=blkPendEdge] at (tap_da1) {};

% Canales horizontales limpios fuera de zSens (CERO solapamiento con etiquetas de zona)
\coordinate (rail_dt1_y) at ([yshift=0.34cm]zSens.north);
\coordinate (rail_t_y)   at ([yshift=0.84cm]zSens.north);
\coordinate (rail_da1_y) at ([yshift=-0.34cm]zSens.south);
\coordinate (rail_a_y)   at ([yshift=-0.80cm]zSens.south);

% Puertos de entrada separados en cada controlador (posiciones Y exactas)
% BALANCE_PID (5 puertos: de abajo hacia arriba dt2, da2, t, da1, a_tilde)
\coordinate (bal_p1) at ([yshift=-2*\cbPortGap]modBal.west);
\coordinate (bal_p2) at ([yshift=-1*\cbPortGap]modBal.west);
\coordinate (bal_p3) at (modBal.west);
\coordinate (bal_p4) at ([yshift= 1*\cbPortGap]modBal.west);
\coordinate (bal_p5) at ([yshift= 2*\cbPortGap]modBal.west);

% SWING_UP (4 puertos: dt1, t, da1, alpha/a_tilde)
\coordinate (sw_p1) at ([yshift=-1.5*\cbPortGap]modSw.west);
\coordinate (sw_p2) at ([yshift=-0.5*\cbPortGap]modSw.west);
\coordinate (sw_p3) at ([yshift= 0.5*\cbPortGap]modSw.west);
\coordinate (sw_p4) at ([yshift= 1.5*\cbPortGap]modSw.west);

% RECOVER_HOME (3 puertos: t, da1, alpha)
\coordinate (rec_p1) at ([yshift=-1*\cbPortGap]modRec.west);
\coordinate (rec_p2) at (modRec.west);
\coordinate (rec_p3) at ([yshift= 1*\cbPortGap]modRec.west);

% Servo 0° (2 puertos: dt1, t)
\coordinate (pos_p1) at ([yshift=-0.55*\cbPortGap]modPos.west);
\coordinate (pos_p2) at ([yshift= 0.55*\cbPortGap]modPos.west);

% MOTOR_PID_VEL (1 puerto: dt1)
\coordinate (vel_p1) at (modVel.west);

% Destinos separados de guardas en Supervisor FSM (estrictamente a la IZQUIERDA de muxTop -> CERO cruces con sel/reset)
\coordinate (sup_g_a)   at ([xshift=0.45cm]zSup.south west);
\coordinate (sup_g_da1) at ([xshift=0.95cm]zSup.south west);
\coordinate (sup_g_t)   at ([xshift=1.45cm]zSup.south west);
\coordinate (corr_t_y)   at ([yshift=0.42cm]zBank.north);
\coordinate (corr_da1_y) at ([yshift=0.88cm]zBank.north);
\coordinate (corr_a_y)   at ([yshift=1.34cm]zBank.north);

% --- 6 pistas explícitas (posición separada de velocidad, brazo separado de péndulo) ---
% 1) PISTA 1: dtheta_f2 (Teal) -> EXCLUSIVA BALANCE_PID (bal_p1)
\draw[fbT, rounded corners=3pt]
    (lpf2_t.west) -- (trk_dt2 |- lpf2_t.west) -- (trk_dt2 |- bal_p1) -- (bal_p1);
\node[wirelbl, text=blkArmEdge, anchor=south]
    at ([xshift=-0.72cm, yshift=1.5pt]lpf2_t.west) {$\dot{\theta}_{f2}$};
\node[wirelbl, font=\scriptsize\sffamily, text=blkArmEdge, anchor=north]
    at ($(trk_dt2 |- bal_p1)!0.52!(zBank.west |- bal_p1) + (0, -0.5pt)$) {$\dot{\theta}_{f2}$};

% 2) PISTA 2: dalpha_f2 (Navy) -> EXCLUSIVA BALANCE_PID (bal_p2) [salta sobre lpf2_t horizontal]
\draw[lineA, rounded corners=3pt]
    (lpf2_a.west) -- (trk_da2 |- lpf2_a.west)
    -- ($(trk_da2 |- lpf2_t.west)+(0, -\cbHopRad)$)
    arc[start angle=-90, end angle=90, radius=\cbHopRad];
\draw[fbA, rounded corners=3pt]
    ($(trk_da2 |- lpf2_t.west)+(0, \cbHopRad)$) -- (trk_da2 |- bal_p2) -- (bal_p2);
\node[wirelbl, text=blkPendEdge, anchor=north]
    at ([xshift=-0.72cm, yshift=-1.5pt]lpf2_a.west) {$\dot{\alpha}_{f2}$};
\node[wirelbl, font=\scriptsize\sffamily, text=blkPendEdge, anchor=south]
    at ($(trk_dt2 |- bal_p2)!0.52!(zBank.west |- bal_p2) + (0, 0.5pt)$) {$\dot{\alpha}_{f2}$};

% 3) PISTA 3: dtheta_f1 (Teal) -> SWING_UP, Servo 0°, MOTOR_PID_VEL [salta trk_dt2 y trk_da2 en rail_dt1_y]
\draw[lineT, rounded corners=3pt]
    (tap_dt1) -- (tap_dt1 |- rail_dt1_y)
    -- ($(trk_dt2 |- rail_dt1_y)+(\cbHopRad,0)$)
    arc[start angle=0, end angle=180, radius=\cbHopRad]
    -- ($(trk_da2 |- rail_dt1_y)+(\cbHopRad,0)$)
    arc[start angle=0, end angle=180, radius=\cbHopRad];
\draw[lineT, rounded corners=3pt]
    ($(trk_da2 |- rail_dt1_y)+(-\cbHopRad,0)$) -- (trk_dt1 |- rail_dt1_y) -- (trk_dt1 |- vel_p1);
\node[wirelbl, text=blkArmEdge, anchor=south]
    at ($(lpf2_t.north |- rail_dt1_y) + (-0.25cm, 1pt)$) {$\dot{\theta}_{f1}$};

\node[dot=blkArmEdge] at (trk_dt1 |- sw_p1) {};
\draw[fbT] (trk_dt1 |- sw_p1) -- (sw_p1);
\node[wirelbl, font=\scriptsize\sffamily, text=blkArmEdge, anchor=north]
    at ($(trk_dt2 |- sw_p1)!0.52!(zBank.west |- sw_p1) + (0, -0.5pt)$) {$\dot{\theta}_{f1}$};

\node[dot=blkArmEdge] at (trk_dt1 |- pos_p1) {};
\draw[fbT] (trk_dt1 |- pos_p1) -- (pos_p1);
\node[wirelbl, font=\scriptsize\sffamily, text=blkArmEdge, anchor=north]
    at ($(trk_dt2 |- pos_p1)!0.52!(zBank.west |- pos_p1) + (0, -0.5pt)$) {$\dot{\theta}_{f1}$};

\draw[fbT] (trk_dt1 |- vel_p1) -- (vel_p1);
\node[wirelbl, font=\scriptsize\sffamily, text=blkArmEdge, anchor=south]
    at ($(trk_dt2 |- vel_p1)!0.52!(zBank.west |- vel_p1) + (0, 1pt)$) {$\dot{\theta}_{f1}$};

% 4) PISTA 4: theta (Teal) -> BALANCE, SWING, RECOVER, Servo 0° y Supervisor FSM
\draw[lineT, rounded corners=3pt]
    (tap_t) -- (tap_t |- rail_t_y)
    -- ($(trk_dt2 |- rail_t_y)+(\cbHopRad,0)$) arc[start angle=0, end angle=180, radius=\cbHopRad]
    -- ($(trk_da2 |- rail_t_y)+(\cbHopRad,0)$) arc[start angle=0, end angle=180, radius=\cbHopRad]
    -- ($(trk_dt1 |- rail_t_y)+(\cbHopRad,0)$) arc[start angle=0, end angle=180, radius=\cbHopRad];
\draw[fbT, rounded corners=3pt]
    ($(trk_dt1 |- rail_t_y)+(-\cbHopRad,0)$) -- (trk_t |- rail_t_y)
    -- (trk_t |- corr_t_y) -| (sup_g_t);
\node[wirelbl, text=blkArmEdge, anchor=south]
    at ($(lpf1_t.north |- rail_t_y) + (-0.15cm, 1pt)$) {$\theta$};
\node[wirelbl, text=blkArmEdge, anchor=south]
    at ($(zBank.north west |- corr_t_y) + (-0.35cm, 1pt)$) {$\theta$};

% Taps de theta hacia BALANCE (bal_p3), SWING (sw_p2), RECOVER (rec_p1), Servo 0° (pos_p2)
\foreach \pt in {bal_p3, sw_p2, rec_p1, pos_p2} {
    \node[dot=blkArmEdge] at (trk_t |- \pt) {};
    \draw[lineT] (trk_t |- \pt) -- ($(trk_dt1 |- \pt)+(-\cbHopRad,0)$)
        arc[start angle=180, end angle=0, radius=\cbHopRad];
    \draw[fbT] ($(trk_dt1 |- \pt)+(\cbHopRad,0)$) -- (\pt);
    \node[wirelbl, font=\scriptsize\sffamily, text=blkArmEdge, anchor=south]
        at ($(trk_dt2 |- \pt)!0.52!(zBank.west |- \pt) + (0, 0.5pt)$) {$\theta$};
}

% 5) PISTA 5: dalpha_f1 (Navy) -> BALANCE, SWING, RECOVER y Supervisor FSM
\draw[fbA, rounded corners=3pt]
    (tap_da1) -- (tap_da1 |- rail_da1_y)
    -- (trk_da1 |- rail_da1_y)
    -- (trk_da1 |- corr_da1_y) -| (sup_g_da1);
\node[wirelbl, text=blkPendEdge, anchor=north]
    at ($(lpf2_a.south |- rail_da1_y) + (-0.25cm, -1pt)$) {$\dot{\alpha}_{f1}$};
\node[wirelbl, text=blkPendEdge, anchor=south]
    at ($(zBank.north west |- corr_da1_y) + (-0.35cm, 1pt)$) {$\dot{\alpha}_{f1}$};

\foreach \pt in {bal_p4, sw_p3, rec_p2} {
    \node[dot=blkPendEdge] at (trk_da1 |- \pt) {};
    \draw[lineA] (trk_da1 |- \pt)
        -- ($(trk_t |- \pt)+(-\cbHopRad,0)$) arc[start angle=180, end angle=0, radius=\cbHopRad]
        -- ($(trk_dt1 |- \pt)+(-\cbHopRad,0)$) arc[start angle=180, end angle=0, radius=\cbHopRad];
    \draw[fbA] ($(trk_dt1 |- \pt)+(\cbHopRad,0)$) -- (\pt);
    \node[wirelbl, font=\scriptsize\sffamily, text=blkPendEdge, anchor=south]
        at ($(trk_dt2 |- \pt)!0.52!(zBank.west |- \pt) + (0, 0.5pt)$) {$\dot{\alpha}_{f1}$};
}

% 6) PISTA 6: alpha, alpha_tilde (Navy) -> BALANCE, SWING, RECOVER y Supervisor FSM
\draw[fbA, rounded corners=3pt]
    (tap_a) -- (tap_a |- rail_a_y)
    -- (trk_a |- rail_a_y)
    -- (trk_a |- corr_a_y) -| (sup_g_a);
\node[wirelbl, text=blkPendEdge, anchor=north]
    at ($(deriv_a.south |- rail_a_y) + (-0.25cm, -1pt)$) {$\alpha,\,\tilde{\alpha}$};
\node[wirelbl, text=blkPendEdge, anchor=south]
    at ($(zBank.north west |- corr_a_y) + (-0.35cm, 1pt)$) {$\alpha,\,\tilde{\alpha}$};

\foreach \pt/\lb in {bal_p5/{$\tilde{\alpha}$}, sw_p4/{$\alpha,\tilde{\alpha}$}, rec_p3/{$\alpha$}} {
    \node[dot=blkPendEdge] at (trk_a |- \pt) {};
    \draw[lineA] (trk_a |- \pt)
        -- ($(trk_da1 |- \pt)+(-\cbHopRad,0)$) arc[start angle=180, end angle=0, radius=\cbHopRad]
        -- ($(trk_t |- \pt)+(-\cbHopRad,0)$) arc[start angle=180, end angle=0, radius=\cbHopRad]
        -- ($(trk_dt1 |- \pt)+(-\cbHopRad,0)$) arc[start angle=180, end angle=0, radius=\cbHopRad];
    \draw[fbA] ($(trk_dt1 |- \pt)+(\cbHopRad,0)$) -- (\pt);
    \node[wirelbl, font=\scriptsize\sffamily, text=blkPendEdge, anchor=south]
        at ($(trk_dt2 |- \pt)!0.52!(zBank.west |- \pt) + (0, 0.5pt)$) {\lb};
}

% --- Selector FSM -> MUX FSM ---
\draw[adapt] (muxTop |- zSup.south) -- node[wirelbl, right=4pt, text=blkDistEdge]{\textbf{sel / reset}} (muxTop);

\end{tikzpicture}%
}
\caption{Arquitectura global de conmutaci\'on e interconexi\'on de los controladores.}
\label{fig:conmutacion_bloques}
\end{figure*}
